\section{Conclusion}\label{sec:conclusion}
In this demo paper, we present a system for online schema alignment.
Through this demonstration, we show that, in unindexed and decentralized networks, \gls{ltqp} queries can be executed successfully in a generalized manner without increasing the complexity of the user-issued queries.
Limitation remain in the expressivity of those rules.
Future work includes extending the system to support more advanced reasoning mechanisms, such as the SPARQL entailment regimes, full \gls{sssom}, Notation3 and the upcoming SHACL rules.
%Engines such as the EYE reasoner or modern Prolog implementations like Scryer Prolog could enable such reasoning capabilities.

%\rt{What is missing for the in this conclusion, is a clear view on the limitations of the current work, and an idea on how these can be solved.}
%\rt{Also, make sure to end with some high-level perspectives. Come back to the need of this work. What is the impact? What does it mean for KG tools and applications? What is the impact on the broader society?}
